% Historical compatibility filename; this is engineering rationale, not a proof.
\documentclass[12pt]{article}
\usepackage[margin=1in]{geometry}
\usepackage{amsmath}
\usepackage[colorlinks=true,urlcolor=blue,linkcolor=blue]{hyperref}

\title{BIOS Engineering Rationale and Evidence Boundaries}
\author{SIH Fleet Priority Team}
\date{September 6, 2026}

\begin{document}
\maketitle

\section*{Scope and current status}
This document retains its historical filename for compatibility. It replaces
unsupported mathematical and percentage claims with implemented mechanisms,
conditional assumptions and reproducible evidence. BIOS 7 is a candidate under
verification, not a released default.

The frozen \texttt{5cce5de} fixed-floor, 50-AMR, seed-0 run ended at 1,200 simulated
seconds with 99/100 tasks for both current-source BIOS 6 and BIOS 7, and zero
observed robot--robot, robot--human and robot--rack contacts. This is a liveness
failure, not a release pass. Later cache, charger and idle-tail changes still need
a fresh passing campaign; no pending result is assumed successful. See the
\href{../artifacts/benchmarks/bios7-release-5cce5de-regression50.json}{retained raw regression}
and \href{24-BIOS7-RELEASE-CHECKLIST.md}{release checklist}.

\section{Separate allocation, traffic and protective stopping}
Task allocation chooses work, traffic coordination negotiates shared space, and
local stopping reacts to sensors. The stop-and-wait reference can use the same
task allocator as BIOS. It must not be described as necessarily lacking
communication, energy checks or multiple tasks. It waits while movement is blocked
and may resume when space clears; persistent circular waits are possible, but not
every stop is permanent.

Centralized references are separate comparisons. A failed repository reference
does not establish that every centralized planner or industrial fleet is inferior.
Known reference limitations are recorded in the
\href{23-BIOS7-EXPERIMENT.md}{BIOS 7 experiment notes}.

\section{Why the mechanisms may help}
\begin{itemize}
\item \textbf{Priorities and expiring leases.} Deterministic tie-breaking, bounded
reservations and local recovery help resolve competing requests. Lease expiry
removes stale software authority, not a failed robot's physical body. Progress
still requires an available safe route or yielding space.
\item \textbf{Energy-aware allocation.} Estimates include approach, cargo-dependent
travel and handling, and return to a configured reachable charger. Battery reserve
and hard deadlines are checked against these estimates. Real battery consumption
and completion time can differ from the model. No-dock legacy maps do not prove
return-to-dock feasibility.
\item \textbf{BIOS 6 and Auction V2.} Event-triggered messages, congestion experience
and bounded future reservations aim to reduce redundant communication and waiting.
A bounded future slot does not bound every candidate evaluation; auction opening
can still inspect the available task catalog.
\item \textbf{BIOS 7 candidate.} Observing a loaded robot complete a corridor
passage can release that task's auction-admission pressure. Physical occupancy,
movement leases and local stopping remain authoritative. Stale or missing
observations retain conservative admission.
\end{itemize}

\section{Conditional assumptions, not universal guarantees}
Progress requires a feasible task catalog, reachable destinations, sufficient
energy and service capacity, space to yield, accurate enough sensing/localization,
functioning command/watchdog paths and the selected protocol's communication
conditions. Permanent obstruction, saturation or failures can prevent completion.
Graph-level reasoning alone does not establish safety for continuous chassis motion.

No proof here establishes universal deadlock freedom, a fixed time-to-deadlock
complexity, logarithmic or linear fleet speedup, or a formula connecting prediction
accuracy directly to collision probability. Communication costs computation and
bandwidth, and predictions can be wrong.

The \href{https://groups.csail.mit.edu/tds/papers/Lynch/jacm85.pdf}{Fischer--Lynch--Paterson result}
rules out guaranteed termination of deterministic consensus in every fully
asynchronous execution with a possible crash, even with reliable messages. It
does not mean agreement is impossible in all executions or validate this robot
implementation.

\section{How evidence must be read}
The \href{../artifacts/benchmarks/bios7-fixed-control-final-development.json}{historical development comparison}
used candidate snapshot \texttt{12d1ac2} and untouched control \texttt{8e4ab51}.
Each BIOS policy completed 90 tasks across three ten-AMR runs with zero observed
contacts. These are finite historical observations, not current-release approval
or 90 independent runs. Different campaign sizes are not a success-rate or speedup
gain.

For safely completed runs with identical inputs, completion-time reduction is
\[
100\left(1-\frac{T_{\mathrm{candidate}}}{T_{\mathrm{baseline}}}\right)\%.
\]
If the baseline reaches a simulation cutoff unfinished, use its observed window
only for a conservative lower bound, not an exact completion-time comparison.
An external worker error or resource timeout gives no performance score.
A contact-free prefix does not establish future safety.

Report completion, all contact categories, messages, bytes, computation, scenario
coverage and determinism separately. Shared-host headless timings are not unloaded
live-control or Raspberry Pi measurements. Retain failed and unrun cases, including
the \href{../artifacts/benchmarks/bios7-capacity-screen.json}{earlier 50-AMR rack-contact failure}.

\section{Promotion and physical deployment}
The \href{24-BIOS7-RELEASE-CHECKLIST.md}{release checklist} requires the fixed
failed-case retest, registered fleets/workloads, SIH overlap comparisons, measured
human/fault interactions, strict same-source V6 nonregression, semantic repeats and
independent-controller evidence. Previously observed seeds must be identified as
retests. A gain in one metric cannot hide failed safety or completion.

The local stop logic is software, not a certified protective device. A real AMR
requires a supported driver, validated sensing/localization and braking,
commissioning, and an independent manufacturer/integrator safety chain. No
physical safety certification, universal hardware compatibility or universal
completion claim is made.

\end{document}
